\documentclass[12pt]{article}
\usepackage[letterpaper,margin=1in]{geometry}
\usepackage[T1]{fontenc}
\usepackage{amsmath,amsthm,mathtools,amssymb}
\usepackage{microtype}
\usepackage{booktabs,longtable,tabularx,graphicx,placeins,setspace}
\usepackage{enumitem}
\usepackage{tikz}
\usetikzlibrary{arrows.meta}
\tikzset{
  dag/.style={
    x=1.10cm,y=.82cm,
    font=\small,
    line width=.4pt,
    >={Triangle[length=3.2mm,width=1.7mm]},
    every label/.style={inner sep=1.5pt}
  },
  dag node/.style={
    circle,draw=black,minimum size=4.8pt,
    inner sep=0pt,outer sep=0pt
  },
  dag observed/.style={dag node,fill=black},
  dag latent/.style={dag node,fill=white},
  dag arrow/.style={->,shorten >=0pt,shorten <=0pt}
}
\usepackage{natbib}
\setcitestyle{aysep={}}
\setlength{\bibsep}{2pt}
\AtBeginDocument{\bibliographystyle{apsr-natbib}}

\usepackage{xcolor}
\definecolor{DarkBlue}{HTML}{00008B}
\usepackage{caption}
\captionsetup{labelfont=bf}
\usepackage[colorlinks=true,allcolors=DarkBlue]{hyperref}
\usepackage{titlesec}
\allowdisplaybreaks
\sloppy

\newcommand{\E}{\mathbb{E}}
\newcommand{\PP}{\mathbb{P}}
\newcommand{\indep}{\mathrel{\perp\!\!\!\perp}}
\DeclareMathOperator{\Var}{Var}
\DeclareMathOperator{\Cov}{Cov}
\DeclareMathOperator{\Corr}{Corr}
\DeclareMathOperator{\SD}{SD}

\theoremstyle{plain}
\newtheorem{proposition}{Proposition}
\newtheorem{corollary}{Corollary}
\theoremstyle{definition}
\newtheorem{assumption}{Assumption}
\newtheorem{definition}{Definition}
\theoremstyle{remark}
\newtheorem*{remark}{Remark}


\title{Manuscript Title}
\author{Author One\thanks{Institution One}\and Author Two\thanks{Institution Two}}
\date{\today}

\makeatletter
\let\SIoriginalmaketitle\@maketitle
\let\SIoriginaland\and
\makeatother
\begin{document}
\maketitle

\begin{abstract}
\noindent Replace this paragraph with the research question, approach, main finding, and substantive implication. The examples below illustrate the manuscript's formatting and cross-references.
\end{abstract}

\noindent\textbf{\em Keywords:} keyword one; keyword two; keyword three

\clearpage

\section{Introduction}\label{sec:intro}

The sample mean summarizes a set of observations with a single quantity. This example separates its definition from the assumptions that establish its sampling properties. Example references include a book on experimental design \citep{fisher1935design} and the methodological article by \citet{athey2019generalized}.

Section~\ref{sec:framework} states the result. SI Section~\ref{SI-sec:proof} provides its proof, and SI Section~\ref{SI-sec:extension} gives a short extension.

\section{Framework}\label{sec:framework}

\begin{assumption}[Independent observations]\label{ass:sampling}
For $i=1,\ldots,n$, with $n\geq 1$:
\begin{enumerate}[label=(\alph*)]
\item $Y_1,\ldots,Y_n$ are mutually independent.
\item $\E[Y_i]=\mu$ and $\Var(Y_i)=\sigma^2<\infty$.
\end{enumerate}
\end{assumption}

\begin{definition}[Sample mean]\label{def:mean}
The {\em sample mean} is
\begin{equation}\label{eq:mean}
\bar Y=\frac{1}{n}\sum_{i=1}^{n}Y_i.
\end{equation}
\end{definition}

\begin{proposition}[Mean and variance]\label{prop:mean}
Under Assumption~\ref{ass:sampling}, the estimator in Definition~\ref{def:mean} satisfies
\begin{equation}\label{eq:moments}
\E[\bar Y]=\mu,
\qquad
\Var(\bar Y)=\frac{\sigma^2}{n}.
\end{equation}
\end{proposition}

The proof appears in SI Section~\ref{SI-sec:proof}. Equation~\eqref{eq:moments} shows how the variance depends on the number of observations.

\clearpage
\section{Illustration and discussion}\label{sec:illustration}

Table~\ref{tab:notation} collects the quantities in Equation~\eqref{eq:mean}. Figure~\ref{fig:diagram} shows an editable diagram. These examples contain no empirical results.

\begin{table}[!htbp]
\centering
\caption{Illustrative quantities.}
\label{tab:notation}
\begin{tabularx}{\linewidth}{@{}lX@{}}
\toprule
Symbol & Meaning \\
\midrule
$Y_i$ & Observation $i$ \\
$\mu$ & Common population mean \\
$\bar Y$ & Sample mean \\
\bottomrule
\end{tabularx}
\end{table}

\begin{figure}[!htbp]
\centering
\begin{tikzpicture}[dag]
\node[dag observed,label=above:{$X$}] (X) at (0,0) {};
\node[dag observed,label=above:{$Y$}] (Y) at (3,0) {};
\draw[dag arrow] (X) -- (Y);
\end{tikzpicture}

\caption{Illustrative directed graph. The arrow represents a proposed effect of $X$ on $Y$.}
\label{fig:diagram}
\end{figure}

The SI also illustrates a weighted mean in Definition~\ref{SI-def:weighted}, with moments in Equation~\eqref{SI-eq:weighted-moments}. Table~\ref{SI-tab:weights} and Figure~\ref{SI-fig:weights} show its construction. The discussion of a new paper should connect its findings to the substantive question and identify the limits of the evidence.

\FloatBarrier
\ifdefined\printbibliography
\printbibliography[title={References}]
\else
\bibliography{references}
\fi

\clearpage
\setcounter{page}{1}
\renewcommand{\thepage}{S\arabic{page}}
\setcounter{section}{0}
\setcounter{equation}{0}
\setcounter{table}{0}
\setcounter{figure}{0}
\setcounter{proposition}{0}
\setcounter{corollary}{0}
\setcounter{assumption}{0}
\setcounter{definition}{0}
\renewcommand{\thesection}{S\arabic{section}}
\renewcommand{\theequation}{S\arabic{equation}}
\renewcommand{\thetable}{S\arabic{table}}
\renewcommand{\thefigure}{S\arabic{figure}}
\renewcommand{\theproposition}{S\arabic{proposition}}
\renewcommand{\thecorollary}{S\arabic{corollary}}
\renewcommand{\theassumption}{S\arabic{assumption}}
\renewcommand{\thedefinition}{S\arabic{definition}}
\renewcommand{\theHsection}{S\arabic{section}}
\renewcommand{\theHequation}{S\arabic{equation}}
\renewcommand{\theHtable}{S\arabic{table}}
\renewcommand{\theHfigure}{S\arabic{figure}}
\renewcommand{\theHproposition}{S\arabic{proposition}}
\providecommand{\theHcorollary}{\arabic{corollary}}
\providecommand{\theHassumption}{\arabic{assumption}}
\providecommand{\theHdefinition}{\arabic{definition}}
\renewcommand{\theHcorollary}{S\arabic{corollary}}
\renewcommand{\theHassumption}{S\arabic{assumption}}
\renewcommand{\theHdefinition}{S\arabic{definition}}
\allowdisplaybreaks[0]
\makeatletter
\let\SIoriginaladdtocontents\addtocontents
\renewcommand{\addtocontents}[2]{
  \def\SIextension{#1}
  \def\SItoc{toc}
  \ifx\SIextension\SItoc
    \SIoriginaladdtocontents{sitoc}{#2}
  \else
    \SIoriginaladdtocontents{#1}{#2}
  \fi
}
\renewcommand*{\l@section}[2]{\begingroup\bfseries\@dottedtocline{1}{0em}{2.2em}{#1}{#2}\endgroup}

\begingroup
\def\@title{Supplementary Information for\\``Manuscript Title''}
\def\@author{Author One \and Author Two}
\def\@date{}
\let\and\SIoriginaland
\SIoriginalmaketitle
\endgroup
\thispagestyle{plain}
\makeatother
{\setcounter{tocdepth}{1}\makeatletter
\let\SIoriginalstarttoc\@starttoc
\renewcommand{\@starttoc}[1]{\SIoriginalstarttoc{sitoc}}
\tableofcontents
\makeatother}
\clearpage
\onehalfspacing

\section{Proof of the main result}\label{SI-sec:proof}

This section proves Proposition~\ref{prop:mean} from Section~\ref{sec:framework} of the main text.

\begin{proof}[Proof of Proposition~\ref{prop:mean}]
By Assumption~\ref{ass:sampling} and Equation~\eqref{eq:mean},
\begin{align}
\E[\bar Y]
&=\frac{1}{n}\sum_{i=1}^{n}\E[Y_i]
=\mu,\label{SI-eq:expectation}\\
\Var(\bar Y)
&=\frac{1}{n^2}\sum_{i=1}^{n}\Var(Y_i)
=\frac{\sigma^2}{n}.\label{SI-eq:variance}
\end{align}
Independence makes the covariance terms zero. These equalities establish Equation~\eqref{eq:moments}.
\end{proof}

\section{A weighted extension}\label{SI-sec:extension}

\subsection{Conditions and result}\label{SI-sec:weighted-result}

\begin{assumption}[Fixed weights]\label{SI-ass:weights}
The nonrandom weights $w_1,\ldots,w_n$ are nonnegative and satisfy $\sum_{i=1}^{n}w_i=1$.
\end{assumption}

\begin{definition}[Weighted mean]\label{SI-def:weighted}
The weighted mean is $\sum_{i=1}^{n}w_iY_i$.
\end{definition}

\begin{proposition}[Weighted moments]\label{SI-prop:weighted}
Under Assumptions~\ref{ass:sampling} and~\ref{SI-ass:weights},
\begin{equation}\label{SI-eq:weighted-moments}
\E\left[\sum_{i=1}^{n}w_iY_i\right]=\mu,
\qquad
\Var\left(\sum_{i=1}^{n}w_iY_i\right)=\sigma^2\sum_{i=1}^{n}w_i^2.
\end{equation}
\end{proposition}

\begin{proof}
Linearity of expectation and $\sum_i w_i=1$ give the mean. Independence gives the variance by the same argument as in Equation~\eqref{SI-eq:variance}.
\end{proof}

\begin{corollary}[Equal weights]\label{SI-cor:equal}
Setting $w_i=1/n$ in Proposition~\ref{SI-prop:weighted} recovers Proposition~\ref{prop:mean}.
\end{corollary}

\clearpage
\subsection{Illustration}\label{SI-sec:illustration}

Table~\ref{SI-tab:weights} gives two weight choices. Figure~\ref{SI-fig:weights} illustrates the construction for two observations, using the same shared drawing styles as Figure~\ref{fig:diagram} in the main text.

\begin{table}[!htbp]
\centering
\caption{Illustrative fixed weights.}
\label{SI-tab:weights}
\begin{tabularx}{\linewidth}{@{}Xll@{}}
\toprule
Choice & Weights & Variance \\
\midrule
Equal weights & $w_i=1/n$ & $\sigma^2/n$ \\
First observation only & $w_1=1$, $w_i=0$ for $i>1$ & $\sigma^2$ \\
\bottomrule
\end{tabularx}
\end{table}

\begin{figure}[!htbp]
\centering
\begin{tikzpicture}[dag]
\node[dag observed,label=left:{$Y_1$}] (Y1) at (0,1) {};
\node[dag observed,label=left:{$Y_2$}] (Y2) at (0,-1) {};
\node[dag observed,label=right:{$w_1Y_1+w_2Y_2$}] (mean) at (3,0) {};
\draw[dag arrow] (Y1) -- (mean);
\draw[dag arrow] (Y2) -- (mean);
\end{tikzpicture}

\caption{Two observations contribute to a weighted mean. Arrows indicate the construction, not causal effects.}
\label{SI-fig:weights}
\end{figure}

\FloatBarrier

\end{document}
